\documentclass[12pt]{article}
\usepackage[utf8]{inputenc}
\usepackage[T1]{fontenc}
\usepackage{hyperref}
\title{Wagaya Language Resources for K-12 Curriculum Development}
\author{Compiled by Claude Sonnet 5.5 (claude-sonnet-5-5) for Andr\'as Kornai}
\date{}
\begin{document}
\maketitle

\section{Introduction}
Wagaya (also Wakaya; ISO 639-3 wga) is an extinct Pama-Nyungan language of the Warluwarric subgroup. Its country lies on the Queensland--Northern Territory border, north-east and east of Tennant Creek. Its closest relatives are Yindjilandji, Bularnu and Warluwarra. There are no known native speakers, so any curriculum would be a heritage and revitalization programme based on historical records.

\section{Primary Educational Resources}
None found in search results.

\section{Linguistic and Reference Materials}
\begin{itemize}
\item Wikipedia article on the Wagaya language. It describes Gavan Breen's 1974 \emph{Wakaya Grammar} and the Eastern and Western dialects. The grammar itself was not located online. \url{https://en.wikipedia.org/wiki/Wagaya_language}
\item ``Wagaya language word list'', compiled by Phillip Rankine in 1985, in the Living Archive of Aboriginal Languages collection on Territory Stories. \url{https://territorystories.nt.gov.au/10070/884443}
\item Grambank entry for Wagaya. \url{https://grambank.clld.org/languages/waga1260}
\end{itemize}

\section{Religious and Community Resources}
No scripture or community organisation was found.

\section{Digital and Audio Resources}
No audio was confirmed in search results. Wikipedia mentions PARADISEC collections of Breen's materials, which include related Yindjilandji items, but no URL for them was seen.

\section{Curriculum Development Context}
The language is poorly documented: one wordlist and one grammar. A heritage programme would need to rely on these records and on consultation with descendant communities. It could also draw on the better-documented related languages.

\end{document}